\begingroup
\titlespacing*{\chapter}
  {0pt}
  {0pt}
  {0pt}

\chapter*{ABSTRAK}

\endgroup

\addcontentsline{toc}{chapter}{ABSTRAK}


% ==================================================
% FORMAT DASAR HALAMAN ABSTRAK
% ==================================================

\begingroup

% Single spacing
\setstretch{1}

% Tidak ada spacing tambahan antar-paragraf
\setlength{\parindent}{0pt}
\setlength{\parskip}{0pt}


% ==================================================
% JARAK:
% ABSTRAK -> JUDUL TESIS = 1 SPASI
% ==================================================

\vspace{14.4pt}


% ==================================================
% JUDUL TESIS
%
% - Center
% - Times
% - 14 pt
% - Bold
% - Italic
% - Kapital
% - Single spacing
% ==================================================

{\centering
\normalfont
\bfseries
\fontsize{14pt}{16.8pt}\selectfont

CODA-EDA: \textit{COMPACT DEPENDENCY-AWARE ESTIMATION OF DISTRIBUTION ALGORITHM} UNTUK \textit{POLICY MINING ATTRIBUTE-BASED ACCESS CONTROL} (ABAC) PADA \textit{EDGE-IOT}

\par
}


% ==================================================
% JARAK:
% JUDUL TESIS -> IDENTITAS = 1 SPASI
% ==================================================

\vspace{14.4pt}


% ==================================================
% IDENTITAS MAHASISWA
%
% - Center
% - Times
% - 12 pt
% - Single spacing
% ==================================================

{\centering
\normalfont
\fontsize{12pt}{14.4pt}\selectfont

Oleh\par

\textbf{Ziadan Qowi}\par
\textbf{NIM: 23524019}\par
\textbf{Program Studi Magister Informatika}\par

}


% ==================================================
% JARAK:
% IDENTITAS -> ISI ABSTRAK = 3 SPASI
%
% 3 x 14.4 pt = 43.2 pt
% ==================================================

\vspace{43.2pt}

\justifying

Kontrol akses berbasis atribut (ABAC) memerlukan penambangan kebijakan (\textit{policy mining}) dari log akses karena penyusunan mahal dan rawan galat.
Pada IoT, penambangan idealnya berjalan langsung pada perangkat \textit{edge} bersumber daya terbatas sekaligus diperbarui secara berkala.
Algoritma yang ada umumnya boros memori atau mengabaikan dependensi antar atribut.
Penelitian ini merancang dan menguji CoDA-EDA (\textit{Compact Dependency-Aware Estimation of Distribution Algorithm}) yang menyatukan pembaruan probabilistik tanpa populasi eksplisit, pemodelan dependensi orde satu melalui pohon Chow-Liu, dan pemisahan pembaruan struktur yang dipicu \textit{drift} dari pembaruan parameter dan memiliki jejak memori O(D).
Dengan paradigma \textit{Design Science Research}, CoDA-EDA diuji terhadap enam \textit{baseline} pada Raspberry Pi 5 atas lima dataset ABAC Lab dan validasi antar-generator \textit{Smart Building IoT} menggunakan uji Friedman dan \textit{post-hoc} Wilcoxon berkoreksi Holm-Bonferroni.
Hasilnya, kualitas mengungguli keluarga \textit{compact} univariat yang menguat pada dimensi tinggi, setara iBOA/BOA, namun tertinggal moderat dari ACO/UBK.
Jejak memorinya datar dengan rasio 1,00 berbanding 6,2-9,2 kali pada algoritma berpopulasi eksplisit.
Pemisahan pembaruan berhasil menghemat waktu saat data stabil, serta strategi \textit{warm-start} 2,0-2,6 kali lebih cepat dengan kualitas non-inferior.
Kontribusinya adalah pertukaran rancangan yang koheren bagi \textit{gateway edge-IoT}.

\vspace{14.4pt}
\textbf{Kata kunci:} ABAC, CoDA-EDA, dependensi, \textit{edge-IoT}, inkremental, memori, penambangan kebijakan.

\endgroup

\newpage

\begingroup
\titlespacing*{\chapter}
  {0pt}
  {0pt}
  {0pt}

% ========================
% ABSTRACT (ENGLISH)
% ========================
\chapter*{\textit{ABSTRACT}}

\endgroup

\addcontentsline{toc}{chapter}{\textit{ABSTRACT}}

\begingroup

% Single spacing
\setstretch{1}

% Tidak ada spacing tambahan antar-paragraf
\setlength{\parindent}{0pt}
\setlength{\parskip}{0pt}


% ==================================================
% JARAK:
% ABSTRAK -> JUDUL TESIS = 1 SPASI
% ==================================================

\vspace{14.4pt}


% ==================================================
% JUDUL TESIS
%
% - Center
% - Times
% - 14 pt
% - Bold
% - Italic
% - Kapital
% - Single spacing
% ==================================================

{\centering
\normalfont
\bfseries
\itshape
\fontsize{14pt}{16.8pt}\selectfont

\textit{CODA-EDA: COMPACT DEPENDENCY-AWARE ESTIMATION OF DISTRIBUTION ALGORITHM FOR POLICY MINING ATTRIBUTE-BASED ACCESS CONTROL (ABAC) IN EDGE-IOT}
\par
}


% ==================================================
% JARAK:
% JUDUL TESIS -> IDENTITAS = 1 SPASI
% ==================================================

\vspace{14.4pt}


% ==================================================
% IDENTITAS MAHASISWA
%
% - Center
% - Times
% - 12 pt
% - Single spacing
% ==================================================

{\centering
\normalfont
\fontsize{12pt}{14.4pt}\selectfont

\textit{By}\par

\textbf{Ziadan Qowi}\par
\textbf{NIM: 23524019}\par
\textbf{\textit{Master's Program in Informatics}}\par

}


% ==================================================
% JARAK:
% IDENTITAS -> ISI ABSTRAK = 3 SPASI
%
% 3 x 14.4 pt = 43.2 pt
% ==================================================

\vspace{43.2pt}

\justifying

\textit{Attribute-Based Access Control (ABAC) requires policy mining from access logs because manual policy construction is costly and error-prone.
In IoT environments, policy mining should ideally be performed directly on resource-constrained edge devices while being periodically updated.
Existing algorithms are generally memory-intensive or fail to account for dependencies among attributes.
This study designs and evaluates CoDA-EDA (Compact Dependency-Aware Estimation of Distribution Algorithm), which integrates population-free probabilistic updates, first-order dependency modeling through a Chow-Liu tree, and the separation of drift-triggered structure updates from parameter updates, while maintaining an O(D) memory footprint.
Following the Design Science Research paradigm, CoDA-EDA is evaluated against six baselines on a Raspberry Pi 5 using five ABAC Lab datasets, together with cross-generator validation on a Smart Building IoT dataset.
Statistical significance is assessed using the Friedman test followed by post-hoc Wilcoxon tests with Holm-Bonferroni correction.
The results show that CoDA-EDA achieves higher solution quality than the univariate compact family, with the advantage becoming more pronounced at higher dimensionalities.
Its performance is comparable to iBOA/BOA, although it remains moderately behind ACO/UBK.
Its memory footprint remains effectively constant, with a ratio of 1.00 compared with 6.2-9.2 times higher memory usage for explicit population-based algorithms.
The separation of update mechanisms successfully reduces computation time when the data distribution is stable, while the warm-start strategy is 2.0-2.6 times faster while maintaining non-inferior solution quality.
The main contribution is a coherent design trade-off for policy mining on edge-IoT gateways.}


\vspace{14.4pt}
\textit{\textbf{Keywords:} ABAC, CoDA-EDA, dependency, edge-IoT, incremental, memory, policy mining}.
\endgroup